% TOP_PROBLEM: {"id": 2100502, "title": "Open Problems in Integrable Systems — Bi-Poisson vector spaces", "category_id": 11, "category": {"id": 11, "name": "dynamical_systems", "display_name": "Dynamical Systems", "description": "Problems about long-term behavior of deterministic systems, Hamiltonian dynamics, and periodic orbits.", "slug": "dynamical-systems", "order_index": 11, "created_at": "2026-07-31T15:26:25.671Z"}, "status": "open"}
% FIRST_POSTED: null
% ENTRY_KIND: statement_only
% Imported from: batch08.tex; UnsolvedMath ID 2100502
% Compile twice: lualatex 2100502.tex
\documentclass[11pt]{article}
\usepackage{fontspec}
\setmainfont{Latin Modern Roman}
\usepackage{amsmath,amssymb,mathtools,unicode-math}
\setmathfont{Latin Modern Math}
\usepackage{xurl,hyperref}
\hypersetup{hidelinks,pdftitle={UnsolvedMath Problem 2100502}}
\newcommand{\sref}[2]{\hyperlink{source:#1:#2}{[S#2]}}
\newcommand{\cataloguescope}[1]{\paragraph{Mathematical question.}#1}
\begin{document}
% UnsolvedMath ID 2100502; source catalogue code AMR-020-0502
\setcounter{section}{2100501}
\section{Open Problems in Integrable Systems — Bi-Poisson vector spaces}\label{2100502}
{\small\sffamily UnsolvedMath ID 2100502 · Dynamical Systems}
\subsection{Definitions and mathematical statement}

\paragraph{Shared notation.}$\mathbb N=\{1,2,\ldots\}$, and $\mathbb Z,\mathbb Q,\mathbb R,\mathbb C$ denote the integers, rationals, reals and complex numbers. Any other conventions are specified locally.


Let $V$ be a finite-dimensional complex vector space with alternating bilinear forms $A,B$, and let $J=\{A+\lambda B:\lambda\in\mathbb C\}\cup\{B\}$. Its generic corank $c$ is the minimum corank of these forms. Define \[\operatorname{Aut}(V,J)=\{T\in GL(V):A(Tu,Tv)=A(u,v),\ B(Tu,Tv)=B(u,v)\ \forall u,v\in V\}.\] Thus the automorphisms preserve each form, not merely the pencil as an unparameterized set. The Grassmannian $\operatorname{Gr}_k(V)$ consists of $k$-dimensional linear subspaces. A subspace $L$ is bi-Lagrangian if $A|_{L\times L}=B|_{L\times L}=0$ and $\dim L=(\dim V+c)/2$, the maximum possible common isotropic dimension. Let $LG(V,J)$ be the closed subvariety of the corresponding Grassmannian defined by these equations.
\paragraph{Statement recorded in the supplied source.}
Give necessary and sufficient conditions on $(A,B)$ for $LG(V,J)$ to be a smooth algebraic variety, and classify its $\operatorname{Aut}(V,J)$-orbits. \sref{2100502}{2}
\subsection{Short English statement}
Determine when the variety of maximal common isotropic subspaces is smooth and how pencil automorphisms partition it into orbits.
\subsection{Sources}
\begin{description}
\item[\hypertarget{source:2100502:1}{S1}] ULAM.AI and the UnsolvedMath Contributors, UnsolvedMath: A Curated Collection of Open Mathematics Problems, record 2100502 (AMR-020-0502). CC BY 4.0. \url{https://huggingface.co/datasets/ulamai/UnsolvedMath}.
\item[\hypertarget{source:2100502:2}{S2}] A. Bolsinov, V. S. Matveev, E. Miranda and S. Tabachnikov, Open Problems, Questions, and Challenges in Finite-Dimensional Integrable Systems (2018), arXiv:1804.03737, Problem 5.3 and bi-Lagrangian definition. \url{https://arxiv.org/abs/1804.03737}.
\end{description}
\label{end:2100502}
\end{document}
